\subsection{Neuron Cellular Automaton}
The script for the implementation of the neuron cellular automaton is \verb+neuron.py+. The code is based on \verb+gameOfLifeInteractiveToric.py+, which was available on the 
Studium page of the course.

Now, we will present how the code works. At first, the number of rows and columns of the grid as well the speed of the simulation are defined. 
Additionally, the rules text is added. This text is displayed when clicking the "rules"-button while running the script. 
Next, we have implemented the new function \verb+count_neighbors_by_state+ that can calculate the number of neighbours in a certain state. 
This function is necessary as we will be working with three different states instead of two.
The function takes the grid and the state that we want to determine the number of neighbours for as arguments. 
The neighbours for each gridpoint are found using \verb+numpy.roll+ to include the periodic boundary conditions. 
As the next step we have updated the \verb+lifeStep+ function. Now, we first create the grids of neighbours in states 1 and 2 and apply the rules afterwards. 
State 0 represents the resting neurons, state 1 the ready neurons and state 2 the firing neurons. The function then returns the updated grid. 

In the constructor of the \verb+GameOfLifeApp+ class the visualisation is defined. This includes the buttons that are displayed and their actions. 
The time step for the continuous run is defined as well. 

The following parameters are fixed in the code: the number of rows and columns is set to 30 each. For the random population of the grid, 
the probability for state 0 is 0.7, for state 1 it is 0.2 and for state 2 it is 0.1. Initially, all cells are in state 0 (resting). Clicking once sets them to
state 1 and clicking another time to state 2. 
